\chapter[Alimentación]{Alimentación: de \SI{24}{\volt} a \SI{3.3}{\volt}, \SI{5}{\volt} y 5V\_A}\label{cap:alim}

\textbf{En este capítulo vas a ver:}
\begin{itemize}
  \item Cuánta energía necesita cada parte de la placa (el presupuesto de consumo).
  \item Cómo se protege la entrada de \SI{24}{\volt}: diodo en serie contra polaridad invertida y TVS contra picos.
  \item Qué es un convertidor \emph{buck}, explicado desde cero, y por qué se usa en lugar de un regulador lineal.
  \item Por qué se eligió el LMR36015 en su variante FPWM de \SI{1}{\mega\hertz}, y cuánto cuesta esa decisión.
  \item El cálculo de cada pieza externa de los bucks: divisor de realimentación, bobina, condensadores.
  \item La comprobación de los tiempos mínimos de encendido y apagado en el peor caso.
  \item El filtro RC que produce el riel analógico \net{5V_A} y por qué existe.
  \item Cómo se alimenta la NUCLEO y en qué orden se enciende todo.
\end{itemize}

\section{Qué hay que alimentar: el presupuesto de consumo}\label{sec:c02-presupuesto}

Antes de elegir una sola pieza de la fuente hay que saber cuánta corriente pide cada riel. La placa tiene cuatro redes de alimentación:

\begin{itemize}
  \item \net{24V_PROT}: los \SI{24}{\volt} de la fuente de laboratorio, ya protegidos. Alimenta los dos bucks y las dos fuentes de corriente IEPE.
  \item \net{3V3}: \SI{3.3}{\volt} para lo digital y los convertidores: la parte digital del ADC (VDD del PCM1808), los dos DAC MCP4725, el oscilador y las resistencias de pull-up del bus I$^2$C.
  \item \net{5V}: la salida del buck de \SI{5}{\volt}. Solo alimenta el filtro RC que produce \net{5V_A}.
  \item \net{5V_A}: \SI{5}{\volt} ``analógicos'' y filtrados: los dos op-amps TLV9152 y la parte analógica del ADC (VCC del PCM1808).
\end{itemize}

La tabla~\ref{tab:c02-presupuesto} sale de la sección~1 de \archivo{CALCULOS.md}. Son \textbf{cotas superiores}: se toma el máximo del datasheet, no el típico, para que la fuente nunca quede corta.

\begin{table}[htbp]
\centering
\caption{Corrientes de diseño por riel (cotas superiores).}
\label{tab:c02-presupuesto}
\begin{tabularx}{\linewidth}{@{}L{0.13\linewidth}Y L{0.14\linewidth}@{}}
\toprule
\textbf{Riel} & \textbf{De dónde sale la corriente} & \textbf{Total}\\
\midrule
\net{5V_A} & PCM1808 ICC máx. \SI{11}{\milli\ampere} \doc{SLES177B p.6} + 2 TLV9152, cada uno con 2 amplificadores a \SI{750}{\micro\ampere} máx. \doc{SBOS986E p.12} & \SI{14.0}{\milli\ampere}\\
\net{5V} & Solo alimenta a \net{5V_A} a través de \net{R33} & \SI{14.0}{\milli\ampere}\\
\net{3V3} & PCM1808 IDD \SI{10.2}{\milli\ampere} (a \SI{96}{\kilo\hertz}, el caso más alto) \doc{SLES177B p.6} + 2 MCP4725 a \SI{400}{\micro\ampere} máx. \doc{DS22039D p.3} + pull-ups I$^2$C + oscilador $\leq$\,\SI{10}{\milli\ampere} \estimado & \SI{22.4}{\milli\ampere}\\
\net{24V_PROT} & Se diseña con el límite de la fuente de banco (\SIrange{30}{50}{\milli\ampere}, \archivo{AGENTS.md} §5). Solo las fuentes IEPE piden como máximo \SI{3.98}{} + \SI{2.51}{\milli\ampere} & \SI{50}{\milli\ampere}\\
\bottomrule
\end{tabularx}
\end{table}

Fíjate en el orden de magnitud: \textbf{decenas de miliamperios}. Esta placa consume muy poco. Eso va a pesar en varias decisiones: los anchos de pista no los fija la corriente sino la fabricación (capítulo~\ref{cap:calculos}), y la eficiencia del buck en vacío importa más que su eficiencia a plena carga.

\begin{calculo}[cuánto le pide la placa a la fuente de \SI{24}{\volt}]
Sumemos la potencia de cada consumo, con las cotas de la tabla:
\begin{align*}
P_{3V3} &= \SI{3.3}{\volt}\cdot\SI{22.4}{\milli\ampere} \approx \SI{74}{\milli\watt}\\
P_{5V} &= \SI{5}{\volt}\cdot\SI{14.0}{\milli\ampere} = \SI{70}{\milli\watt}\\
P_{\text{bucks}} &\leq 2\cdot\SI{0.15}{\watt} = \SI{0.30}{\watt} \quad\text{(pérdidas en vacío, sección~\ref{sec:c02-fpwm})}\\
P_{\text{IEPE}} &\approx (\SI{3.98}{} + \SI{2.51}{\milli\ampere})\cdot\SI{23.5}{\volt} \approx \SI{0.15}{\watt}
\end{align*}
Total: unos \SI{0.6}{\watt}, que a \SI{23.5}{\volt} son unos \SI{25}{\milli\ampere} \estimado. Está por debajo de los \SIrange{30}{50}{\milli\ampere} a los que se limita la fuente de banco. Este número se confirma en el banco con el amperímetro \medir.
\end{calculo}

\section{La entrada: J5, D1 y D2}\label{sec:c02-entrada}

La figura~\ref{fig:c02-entrada} muestra la entrada. Los \SI{24}{\volt} llegan por la bornera \net{J5} (red \net{24V_IN}), pasan por el diodo \net{D1} y desde ahí la red se llama \net{24V_PROT}. Entre \net{24V_PROT} y tierra está el TVS \net{D2}. Este bloque no existía en el Altium del equipo: es el cambio C4 (capítulo~\ref{cap:piezas}).

\begin{figure}[htbp]
\centering
\begin{circuitikz}[scale=0.95, transform shape]
  \draw (0,2) node[left]{\net{J5.1}} to[short,o-] (0.6,2) node[above]{\net{24V_IN}} -- (1.2,2)
        to[sD, l=\net{D1} SS34] (3.6,2) coordinate (p);
  \draw (p) to[short,*-] (6.0,2) coordinate (q) to[short,*-] (8.4,2) node[right]{\net{24V_PROT} $\to$ bucks \net{U6}, \net{U7}};
  \node[above] at (4.8,2) {\net{24V_PROT}};
  \draw (p) to[zDo, l_=\net{D2} SMBJ26A, invert] (3.6,0) node[ground]{};
  \draw (q) -- (6.0,1.2) to[short,-o] (6.0,1.0) node[left]{\net{J1.1}};
  \node[font=\footnotesize, align=left, anchor=west] at (6.2,0.7) {bornera \net{J1} (y \net{J7})\\$\to$ \net{IEPE_FEED} $\to$ fuente IEPE};
  \draw (0,0) node[left]{\net{J5.2}} to[short,o-] (0.8,0) node[ground]{};
\end{circuitikz}
\caption{Entrada de alimentación. \net{D1} en serie bloquea la polaridad invertida; \net{D2} (TVS) recorta los picos sobre \net{24V_PROT}. \net{J1} y \net{J7} son borneras en serie con la alimentación de cada fuente IEPE.}
\label{fig:c02-entrada}
\end{figure}

\subsection{D1: protección contra polaridad invertida}

Si alguien conecta la fuente al revés, todos los integrados recibirían \SI{-24}{\volt}, y eso los mata. \net{D1} está \textbf{en serie}: con la polaridad correcta conduce; con la polaridad invertida queda en inversa y no deja pasar corriente.

Es un SS34, un diodo Schottky de \SI{40}{\volt} y \SI{3}{\ampere} \copia[(design.py)]. Se usa un Schottky porque cae menos tensión en directa que un diodo de silicio normal, unos \SI{0.4}{\volt} \estimado. Con la corriente de la placa, eso es nada:

\begin{calculo}[pérdida en D1]
$P_{D1} = I\cdot V_F = \SI{50}{\milli\ampere}\cdot\SI{0.4}{\volt} = \SI{20}{\milli\watt}$, menos de \SI{1}{\degreeCelsius} de calentamiento (\archivo{CALCULOS.md} §6). La consecuencia práctica es que a los bucks no les llegan \SI{24}{\volt} sino unos \SI{23.5}{\volt}. Por eso ese valor aparece en los cálculos de este capítulo.
\end{calculo}

\begin{decision}[diodo Schottky en serie]
\textbf{Qué se decidió:} un Schottky SS34 en serie con la entrada.\\
\textbf{Alternativas:} un MOSFET de canal P como ``diodo ideal'' (casi sin caída) o un diodo en paralelo con fusible (cortocircuita la fuente al invertirla).\\
\textbf{Por qué:} con \SI{50}{\milli\ampere} la caída del Schottky solo cuesta \SI{20}{\milli\watt}, y es una sola pieza \emph{basic} de JLCPCB (C8678). El MOSFET agrega piezas sin una ventaja que se note a este consumo.\\
\textbf{Fuente:} cambio C4 de \archivo{tools/design.py}; cálculo en \archivo{CALCULOS.md} §6.
\end{decision}

\subsection{D2: el TVS contra picos}

\begin{idea}[qué es un TVS]
Un TVS (\emph{Transient Voltage Suppressor}) es un diodo hecho para absorber picos de tensión. Mientras la tensión está por debajo de su \textbf{tensión de trabajo} $V_{RWM}$ casi no conduce (solo una corriente de fuga pequeña). Si un pico supera su \textbf{tensión de ruptura} $V_{BR}$, conduce mucho y ``se come'' la energía del pico, con la tensión sujeta cerca de su \textbf{tensión de recorte} (\emph{clamping}). Va en paralelo con la línea que protege, de la red a tierra. Es como una válvula de alivio en una tubería: cerrada en condiciones normales, se abre si la presión se dispara.
\end{idea}

El SMBJ26A tiene $V_{RWM}=\SI{26}{\volt}$ y $V_{BR}\geq\SI{28.9}{\volt}$ \copia. ¿Por qué 26?
\begin{itemize}
  \item Tiene que quedar \textbf{por encima} de los \SI{24}{\volt} normales con margen. Si su $V_{RWM}$ fuera menor que la tensión de trabajo, conduciría todo el tiempo y se calentaría.
  \item Tiene que quedar \textbf{por debajo} de lo que aguantan las piezas que protege. Los LMR36015 aceptan hasta \SI{60}{\volt} de entrada \doc{SNVSB49D p.1}, así que el recorte transitorio del TVS (unos \SI{42}{\volt} \copia[(REGLAS.md §5)]) no los pone en riesgo.
\end{itemize}

Como \net{D2} va después de \net{D1}, con la polaridad invertida el TVS no ve nada: \net{D1} ya bloqueó.

\begin{ojo}[valores del TVS sin datasheet a mano]
Los valores del SMBJ26A son los típicos de la serie SMBJ y vienen de una copia, no del datasheet del fabricante exacto (\archivo{CALCULOS.md} §1). En la práctica: no alimentes la placa con más de \SI{26}{\volt} continuos. El TVS también fija una tensión máxima de la que se parte para calcular separaciones entre pistas y tiempos del buck (sección~\ref{sec:c02-tiempos} y capítulo~\ref{cap:calculos}).
\end{ojo}

\subsection{J1 y J7: la habilitación de cada canal}

Cada fuente IEPE toma \net{24V_PROT} a través de una bornera de dos pines, \net{J1} (CH1) o \net{J7} (CH2). El pin~1 es \net{24V_PROT} y el pin~2 es \net{IEPE_FEED} (o \net{IEPE_FEED_SV}), que va a la fuente de corriente. Sirven para dos cosas (cambio C5):
\begin{enumerate}
  \item \textbf{Encender o apagar cada canal}: un cable puente entre los dos tornillos lo habilita.
  \item \textbf{Medir la corriente de la fuente IEPE}: en lugar del puente pones el amperímetro en serie y lees directamente la corriente que va al micrófono (capítulo~\ref{cap:banco}).
\end{enumerate}
El Altium original tenía una bornera aparte (\net{J6}) para alimentar el canal~2 con otra fuente. Se eliminó: los dos canales usan los mismos \SI{24}{\volt}.

\section{Qué es un convertidor buck}\label{sec:c02-buck}

Necesitamos bajar de \SI{23.5}{\volt} a \SI{3.3}{\volt} y a \SI{5}{\volt}. Hay dos formas de hacerlo.

\subsection{Opción 1: un regulador lineal}
Un regulador lineal (el clásico LM317 o un LDO) funciona como una resistencia variable: deja caer la tensión que sobra y la convierte en calor. Toda la corriente de la carga pasa por él con la caída completa.

\begin{calculo}[un regulador lineal de \SI{23.5}{\volt} a \SI{3.3}{\volt}]
$P = (V_{IN}-V_{OUT})\cdot I = (\SI{23.5}{\volt}-\SI{3.3}{\volt})\cdot\SI{22.4}{\milli\ampere} \approx \SI{0.45}{\watt}$.
La eficiencia sería $\eta = V_{OUT}/V_{IN} = 3{,}3/23{,}5 \approx \SI{14}{\percent}$: el 86\,\% de la energía se iría en calor.
\end{calculo}

Medio vatio en un encapsulado pequeño ya es un problema térmico, y la eficiencia es pésima.

\subsection{Opción 2: un convertidor conmutado reductor (buck)}

\begin{idea}[el buck en tres pasos]
Un buck no ``quema'' la tensión que sobra. La \textbf{trocea}:
\begin{enumerate}
  \item Un interruptor (un MOSFET) conecta la entrada a una bobina durante un tiempo $t_{ON}$. La corriente por la bobina sube.
  \item Luego el interruptor se abre y un segundo interruptor (o un diodo) conecta la bobina a tierra durante $t_{OFF}$. La bobina, que no deja cambiar su corriente de golpe, la sigue empujando hacia la carga, y esa corriente baja.
  \item Esto se repite a una frecuencia fija $f_{SW}$ (aquí \SI{1}{\mega\hertz}, un millón de veces por segundo). La bobina y el condensador de salida promedian los pulsos y queda una tensión continua con un pequeño rizado.
\end{enumerate}
La fracción del tiempo que el interruptor de arriba está cerrado se llama \textbf{ciclo de trabajo} $D$. En un buck ideal:
\begin{equation}
V_{OUT} = D\cdot V_{IN} \qquad\Rightarrow\qquad D = \frac{V_{OUT}}{V_{IN}}
\label{eq:c02-duty}
\end{equation}
Como idealmente nada disipa (un interruptor cerrado no tiene caída, uno abierto no tiene corriente), la potencia de entrada es igual a la de salida. Por eso la corriente de entrada es \emph{menor} que la de salida: $I_{IN}\approx D\cdot I_{OUT}$.
\end{idea}

La figura~\ref{fig:c02-buck} muestra el circuito básico y la figura~\ref{fig:c02-il} la corriente de la bobina.

\begin{figure}[htbp]
\centering
\begin{circuitikz}[scale=0.95, transform shape]
  \draw (0,0) to[V, l=$V_{IN}$, invert] (0,3) -- (1,3) to[nos, l=$S_1$ (alto)] (3,3) coordinate (sw);
  \draw (sw) to[nos, l_=$S_2$ (bajo)] (3,0);
  \draw (sw) to[L, l=$L$, i>^=$i_L$] (6,3) coordinate (o) to[short,*-] (8,3) to[R, l=carga] (8,0);
  \draw (o) to[C, l=$C_{OUT}$, *-*] (6,0);
  \draw (0,0) -- (8,0);
  \node[anchor=west, font=\footnotesize] at (3.15,2.55) {nodo SW};
  \draw (6,3) ++(0,0.7) node[font=\footnotesize]{$V_{OUT}$};
\end{circuitikz}
\caption{Buck síncrono simplificado. $S_1$ y $S_2$ se cierran alternadamente. En el LMR36015 los dos interruptores son MOSFET integrados (``síncrono''); en un buck clásico, $S_2$ es un diodo.}
\label{fig:c02-buck}
\end{figure}

\begin{figure}[htbp]
\centering
\begin{tikzpicture}
\begin{axis}[width=0.85\linewidth, height=4.4cm, xlabel={tiempo ($\mu$s)}, ylabel={$i_L$ (mA)},
  xmin=0, xmax=3, ymin=-260, ymax=300, grid=both, grid style={gray!20},
  xticklabel style={/pgf/number format/use comma},
  legend style={font=\footnotesize, at={(0.5,-0.42)}, anchor=north, legend columns=2, draw=none}]
  % D = 0.14 a 1 MHz: sube 0.14 us, baja 0.86 us; dIL = 0.42 A (3.3 V, 24 V, 6.8 uH)
  \addplot[thick, cobre] coordinates {(0,-210) (0.14,210) (1,-210) (1.14,210) (2,-210) (2.14,210) (3,-210)};
  \addlegendentry{FPWM sin carga (promedio 0)\quad}
  \addplot[thick, dashed, verdeplaca] coordinates {(0,-188) (0.14,232) (1,-188) (1.14,232) (2,-188) (2.14,232) (3,-188)};
  \addlegendentry{con \SI{22}{\milli\ampere} de carga}
  \addplot[gray, dotted] coordinates {(0,0) (3,0)};
\end{axis}
\end{tikzpicture}
\caption{Forma de la corriente de la bobina de \SI{3.3}{\volt} con \SI{24}{\volt} de entrada a \SI{1}{\mega\hertz} (rizado $\Delta I_L\approx\SI{0.42}{\ampere}$, sección~\ref{sec:c02-bobina}). Sube durante $t_{ON}\approx\SI{0.14}{\micro\second}$ y baja el resto del período. Como la carga es tan pequeña, el triángulo queda casi centrado en cero: en modo FPWM la corriente se hace negativa en parte del ciclo.}
\label{fig:c02-il}
\end{figure}

\begin{calculo}[ciclo de trabajo de los dos bucks]
Con $V_{IN}=\SI{23.5}{\volt}$ (24 V menos \net{D1}):
\[
D_{3V3} = \frac{3{,}3}{23{,}5}\approx 0{,}14 \qquad D_{5V} = \frac{5}{23{,}5}\approx 0{,}21
\]
A \SI{1}{\mega\hertz} el período es \SI{1}{\micro\second}, así que el interruptor de arriba solo está cerrado \SI{140}{\nano\second} (3V3) o \SI{213}{\nano\second} (5V) en cada ciclo. Son pulsos muy cortos; la sección~\ref{sec:c02-tiempos} comprueba que el integrado los puede hacer.
\end{calculo}

\section{El LMR36015 y el modo FPWM}\label{sec:c02-fpwm}

Los dos bucks usan el mismo integrado, \textbf{LMR36015FBRNXR} de Texas Instruments: \net{U6} produce \net{3V3} y \net{U7} produce \net{5V}. Es un buck síncrono de \SIrange{4.2}{60}{\volt} de entrada y hasta \SI{1.5}{\ampere} de salida \doc{SNVSB49D p.1}, en un encapsulado VQFN de \SI{2}{\milli\meter} $\times$ \SI{3}{\milli\meter} (``HotRod'', RNX).

El sufijo importa. TI vende varias variantes del LMR36015 y esta es la \textbf{FB}: salida ajustable, \SI{1}{\mega\hertz} y \textbf{FPWM} \doc{SNVSB49D p.3}.

\begin{idea}[PFM contra FPWM]
Cuando un buck tiene muy poca carga, como el nuestro, puede funcionar de dos maneras \doc{SNVSB49D p.16}:
\begin{itemize}
  \item \textbf{PFM} (\emph{pulse frequency modulation}): manda ráfagas de pulsos solo cuando la salida baja. Gasta muy poco (el LMR36015 en PFM pide unos \SI{26}{\micro\ampere} sin carga \doc{SNVSB49D p.8}), pero \textbf{la frecuencia de conmutación cambia con la carga} y el rizado de salida es mayor.
  \item \textbf{FPWM} (\emph{forced PWM}): conmuta siempre a la frecuencia fija, aunque no haya carga. La corriente de la bobina sigue siendo un triángulo completo, que en vacío baja por debajo de cero (figura~\ref{fig:c02-il}). La frecuencia es fija, pero se pierde energía moviendo esa corriente de un lado a otro.
\end{itemize}
\end{idea}

¿Por qué importa la frecuencia? Porque esta es una placa de \textbf{audio}. Si un buck en PFM hace ráfagas a unos pocos kilohercios, ese ruido cae justo en la banda que queremos medir (\SI{20}{\hertz} a \SI{20}{\kilo\hertz}), y ningún filtro lo separa de la señal. Con FPWM a \SI{1}{\mega\hertz} el ruido de conmutación queda 50 veces por encima de \SI{20}{\kilo\hertz}, donde los filtros RC lo atenúan fácilmente.

\begin{decision}[LMR36015 en variante FPWM de \SI{1}{\mega\hertz}]
\textbf{Qué se decidió:} mantener el LMR36015\textbf{FB}RNXR (FPWM, \SI{1}{\mega\hertz}) para los dos rieles.\\
\textbf{Alternativas:} la variante PFM (más eficiente en vacío pero con frecuencia variable) o un regulador lineal (sin ruido de conmutación pero con unos \SI{0.45}{\watt} de pérdida, sección~\ref{sec:c02-buck}).\\
\textbf{Por qué:} con frecuencia fija el ruido de la fuente nunca cae en la banda de audio. Además, el integrado sigue en conducción continua (CCM) a frecuencia fija aun sin carga \doc{SNVSB49D p.16, tabla 9-1}, así que no hay modos raros a baja carga. El costo es el consumo en vacío.\\
\textbf{Fuente:} \archivo{REGLAS.md} §3 (LMR36015); plan de ruteo aprobado (decisión~10, ``no propongo cambiarlo''); \archivo{CALCULOS.md} §2.
\end{decision}

\subsection*{El costo del FPWM: pérdidas sin carga}

El datasheet no publica la corriente de entrada en FPWM sin carga. \archivo{CALCULOS.md} §2 la acota así \estimado:
\begin{itemize}
  \item conducción con las resistencias máximas de los MOSFET ($R_{DS(on)}$ de \SI{435}{\milli\ohm} el de arriba y \SI{280}{\milli\ohm} el de abajo \doc{SNVSB49D p.7}): unos \SI{6.5}{\milli\watt};
  \item conmutación con flancos de \SI{10}{\nano\second} (supuesto): unos \SI{42}{\milli\watt};
  \item más el regulador interno y la carga de las compuertas, que no están publicados.
\end{itemize}
La resistencia térmica JEDEC del encapsulado es de \SI{72.5}{\degreeCelsius\per\watt} \doc{SNVSB49D p.6}. Con una cota de \SI{0.15}{\watt} por integrado, se calentaría como mucho
\[
\Delta T = P\cdot R_{\theta JA} = \SI{0.15}{\watt}\cdot\SI{72.5}{\degreeCelsius\per\watt}\approx\SI{11}{\degreeCelsius}
\]
por encima del ambiente. En una placa de 4 capas con plano de tierra debería quedar por debajo de eso. Es una de las primeras cosas que hay que medir en el banco: \medir.

\section{El circuito de cada buck}\label{sec:c02-circuito}

La figura~\ref{fig:c02-u6} muestra \net{U6} con todas sus piezas. \net{U7} es idéntico con sus propios designadores (tabla~\ref{tab:c02-piezas}).

\begin{figure}[htbp]
\centering
\begin{circuitikz}[scale=0.84, transform shape]
  % integrado
  \draw[thick, fill=cobre!5] (0,-0.2) rectangle (3,5.6);
  \node[font=\bfseries] at (1.5,5.2) {\net{U6}};
  \node[font=\footnotesize] at (1.5,4.75) {LMR36015};
  \node[font=\scriptsize, anchor=west] at (0.05,3.3) {VIN 2,10};
  \node[font=\scriptsize, anchor=west] at (0.05,2.3) {EN 9};
  \node[font=\scriptsize, anchor=west] at (0.05,1.3) {VCC 5};
  \node[font=\scriptsize, anchor=west] at (0.05,0.55) {PGND 1,11};
  \node[font=\scriptsize, anchor=west] at (0.05,0.15) {AGND 6};
  \node[font=\scriptsize, anchor=east] at (2.95,3.8) {BOOT 4};
  \node[font=\scriptsize, anchor=east] at (2.95,3.0) {SW 12, NC 3};
  \node[font=\scriptsize, anchor=east] at (2.95,1.3) {FB 7};
  % entrada
  \draw (0,3.3) -- (-1.2,3.3) coordinate (vin) -- (-4.6,3.3) node[left]{\net{24V_PROT}};
  \draw (0,2.3) -- (-1.2,2.3) -- (vin); \fill (vin) circle (1.5pt);
  \draw (-2.4,3.3) to[C, l_=\net{C29}, *-] (-2.4,1.5) node[ground]{};
  \draw (-3.8,3.3) to[C, l_={\net{C27}, \net{C28}}, *-] (-3.8,1.5) node[ground]{};
  \node[font=\scriptsize, align=center] at (-3.8,0.65) {2$\times$\SI{4.7}{\micro\farad}};
  \node[font=\scriptsize, align=center] at (-2.4,0.45) {\SI{220}{\nano\farad}\\(y \net{C30})};
  \draw (0,1.3) -- (-0.7,1.3) to[C, l_=\net{C36}] (-0.7,-0.4) node[ground]{};
  \draw (0,0.35) -- (0.0,0.35);
  \draw (1.5,-0.2) -- (1.5,-0.6) node[ground]{};
  % BOOT y SW
  \draw (3,3.8) -- (4.4,3.8) to[C] (4.4,3.0);
  \node[font=\footnotesize, above] at (4.4,3.85) {\net{C35}};
  \draw (3,3.0) -- (4.4,3.0) to[short, *-] (4.8,3.0) to[L, l=\net{L1} \SI{6.8}{\micro\henry}] (6.6,3.0) coordinate (out);
  % realimentación
  \draw (out) to[short, *-] (7.6,3.0) coordinate (rt);
  \draw (rt) to[R, l=\net{R29} \SI{100}{\kilo\ohm}, *-] (7.6,1.3) coordinate (fb) to[R, l=\net{R30} \SI{43.2}{\kilo\ohm}] (7.6,-0.4) node[ground]{};
  \draw (fb) to[short, *-] (3,1.3);
  \node[font=\scriptsize, above] at (5.2,1.3) {\net{FB_3}};
  % salida
  \draw (rt) -- (10.2,3.0) coordinate (co) to[short, *-o] (11.2,3.0) node[right]{\net{3V3}};
  \draw (co) to[C, l=\net{C31}--\net{C34}] (10.2,0.6) node[ground]{};
  \node[font=\scriptsize, anchor=west] at (10.6,1.25) {4$\times$\SI{15}{\micro\farad}};
\end{circuitikz}
\caption{Buck de \SI{3.3}{\volt} (\net{U6}) con sus piezas, según \archivo{tools/target.json}. EN (pin~9) va a \net{24V_PROT}: el buck está siempre encendido. El pin NC~3 se une a SW (pin~12) en la placa, como pide el datasheet \doc{SNVSB49D p.4}. PG (pin~8) no se usa.}
\label{fig:c02-u6}
\end{figure}

\begin{table}[htbp]
\centering
\caption{Piezas de los dos bucks. Los valores de la columna ``TI'' son los de la tabla de componentes típicos para \SI{1}{\mega\hertz} \doc{SNVSB49D p.19, tabla 10-1}.}
\label{tab:c02-piezas}
\begin{tabularx}{\linewidth}{@{}L{0.19\linewidth}L{0.15\linewidth}L{0.15\linewidth}Y@{}}
\toprule
\textbf{Función} & \textbf{3V3 (\net{U6})} & \textbf{5V (\net{U7})} & \textbf{Valor y comparación con TI}\\
\midrule
CIN (entrada) & \net{C27}, \net{C28} & \net{C37}, \net{C38} & \SI{4.7}{\micro\farad} \SI{100}{\volt} X7R 2220 cada uno. TI: \SI{4.7}{\micro\farad} + 2$\times$\SI{220}{\nano\farad}\\
CHF (alta frec.) & \net{C29}, \net{C30} & \net{C39}, \net{C40} & \SI{220}{\nano\farad} \SI{100}{\volt} X7R 0805, uno por cada par VIN/PGND\\
COUT (salida) & \net{C31}--\net{C34} & \net{C41}--\net{C44} & 4$\times$\SI{15}{\micro\farad} \SI{16}{\volt} X7R 1210. TI: 3$\times$\SI{15}{\micro\farad} (optimizado para transitorios) o 2$\times$ mínimo\\
Bobina & \net{L1} \SI{6.8}{\micro\henry} XAL4030 & \net{L2} \SI{10}{\micro\henry} XAL4040 & TI: \SI{6.8}{\micro\henry} (3,3 V) y \SI{10}{\micro\henry} (5 V)\\
CBOOT & \net{C35} & \net{C45} & \SI{100}{\nano\farad} X7R \SI{50}{\volt}\\
CVCC & \net{C36} & \net{C46} & \SI{2.2}{\micro\farad} X7R 0805\\
RFBT / RFBB & \net{R29} / \net{R30} & \net{R31} / \net{R32} & \SI{100}{\kilo\ohm} / \SI{43.2}{\kilo\ohm} y \SI{100}{\kilo\ohm} / \SI{24.9}{\kilo\ohm}, iguales a los de TI\\
CFF & --- & --- & TI lista \SI{20}{\pico\farad}; esta placa no lo lleva (ver el recuadro ``Ojo'' de la sección~\ref{sec:c02-cond})\\
\bottomrule
\end{tabularx}
\end{table}

Fíjate en que el diseño sigue casi pieza por pieza la tabla de TI para \SI{1}{\mega\hertz}. Eso es deliberado: la tabla del fabricante ya viene con la estabilidad del lazo de control comprobada. En las secciones siguientes se verifica cada valor.

\subsection{El divisor de realimentación: de dónde salen los \SI{3.3}{\volt}}

El buck ajusta su ciclo de trabajo hasta que el pin FB vale exactamente su tensión de referencia interna. Esa referencia está entre \SIrange{0.985}{1.015}{\volt}, con \SI{1}{\volt} típico \doc{SNVSB49D p.6}. Si pones un divisor resistivo entre la salida y tierra con el punto medio en FB, la salida se acomoda para que el punto medio valga \SI{1}{\volt}:

\begin{equation}
V_{OUT} = V_{REF}\left(1 + \frac{R_{FBT}}{R_{FBB}}\right)
\qquad\Longleftrightarrow\qquad
R_{FBB} = \frac{R_{FBT}}{V_{OUT}/V_{REF} - 1}
\label{eq:c02-divisor}
\end{equation}

TI recomienda $R_{FBT}=\SI{100}{\kilo\ohm}$, con un máximo de \SI{1}{\mega\ohm} \doc{SNVSB49D p.20}. Valores más bajos captan menos ruido pero cargan más la salida; más altos ahorran corriente pero hacen el nodo FB más sensible.

\begin{calculo}[divisores de \net{U6} y \net{U7}]
\textbf{3V3:} $R_{FBB}$ ideal $=\dfrac{\SI{100}{\kilo\ohm}}{3{,}3/1-1}=\SI{43.48}{\kilo\ohm}$. El valor comercial E96 más cercano es \SI{43.2}{\kilo\ohm} (\net{R30}). Con él:
\[
V_{OUT} = \SI{1}{\volt}\cdot\left(1+\frac{100}{43{,}2}\right) = \SI{3.315}{\volt}
\]
Con la referencia en sus extremos (\SIrange{0.985}{1.015}{\volt}) la salida queda entre \SIrange{3.265}{3.365}{\volt}, sin contar la tolerancia del \SI{1}{\percent} de las resistencias.

\textbf{5V:} $R_{FBB}$ ideal $=\dfrac{\SI{100}{\kilo\ohm}}{5/1-1}=\SI{25}{\kilo\ohm}$. El E96 más cercano es \SI{24.9}{\kilo\ohm} (\net{R32}), el mismo del ejemplo de \SI{5}{\volt} del datasheet \doc{SNVSB49D p.21}:
\[
V_{OUT} = \SI{1}{\volt}\cdot\left(1+\frac{100}{24{,}9}\right) = \SI{5.016}{\volt}
\]
Entre \SIrange{4.941}{5.091}{\volt} por la tolerancia de la referencia.

\textbf{Corriente del divisor:} $3{,}3/(100+43{,}2)\,\si{\kilo\ohm}\approx\SI{23}{\micro\ampere}$ y $5/(100+24{,}9)\,\si{\kilo\ohm}\approx\SI{40}{\micro\ampere}$. Despreciable frente a los miliamperios de carga.
\end{calculo}

El datasheet garantiza la regulación de la salida en FPWM dentro de $\pm\SI{1.5}{\percent}$ \doc{SNVSB49D p.8}. Ese es el número que se usa en el resto del documento (por ejemplo, para el rango de \net{5V_A} en el capítulo~\ref{cap:analog}).

\subsection{La bobina}\label{sec:c02-bobina}

La bobina se elige por su rizado de corriente $\Delta I_L$, la diferencia entre el pico y el valle del triángulo de la figura~\ref{fig:c02-il}. Durante $t_{ON}$ la bobina tiene $V_{IN}-V_{OUT}$ encima, y como $v=L\,di/dt$:
\begin{equation}
\Delta I_L = \frac{(V_{IN}-V_{OUT})\cdot t_{ON}}{L} = \frac{(V_{IN}-V_{OUT})\cdot V_{OUT}}{V_{IN}\cdot L\cdot f_{SW}}
\label{eq:c02-rizado}
\end{equation}

TI recomienda que el rizado esté entre el 20 y el 40\,\% de la corriente máxima, y aclara que cuando la carga real es mucho menor que la del integrado se use \textbf{la corriente máxima del integrado} (\SI{1.5}{\ampere}) en ese cálculo \doc{SNVSB49D p.21}. Es nuestro caso.

\begin{calculo}[rizado con las bobinas elegidas, a \SI{24}{\volt} y \SI{1}{\mega\hertz}]
\textbf{3V3, \net{L1} = \SI{6.8}{\micro\henry}:}
$\Delta I_L = \dfrac{(24-3{,}3)\cdot 3{,}3}{24\cdot\SI{6.8}{\micro\henry}\cdot\SI{1}{\mega\hertz}} = \SI{0.42}{\ampere}$
$\Rightarrow$ $0{,}42/1{,}5\approx 0{,}28$, el \SI{28}{\percent} de \SI{1.5}{\ampere}.

\textbf{5V, \net{L2} = \SI{10}{\micro\henry}:}
$\Delta I_L = \dfrac{(24-5)\cdot 5}{24\cdot\SI{10}{\micro\henry}\cdot\SI{1}{\mega\hertz}} = \SI{0.40}{\ampere}$
$\Rightarrow$ $0{,}40/1{,}5\approx 0{,}26$, el \SI{26}{\percent}.

Los dos quedan dentro del rango de 20--40\,\% que pide TI.
\end{calculo}

El peor caso es el oscilador en su frecuencia más baja (\SI{0.85}{\mega\hertz}, \doc{SNVSB49D p.7}) con la entrada más alta. Ahí \archivo{CALCULOS.md} §2 da $\Delta I_L=\SI{0.506}{\ampere}$ para 3V3 con \SI{28.9}{\volt}, y la corriente de pico en la bobina llega a unos \SI{0.28}{\ampere}. Con ese pico, y no con los \SI{22}{\milli\ampere} de la carga, se dimensionan las pistas de SW, bobina y COUT (capítulo~\ref{cap:calculos}).

\begin{ojo}[corriente de saturación]
Una bobina se elige también por su corriente de saturación: por encima de ella la inductancia cae y la corriente se dispara. El repositorio no tiene el dato de saturación de la XAL4030-682 ni de la XAL4040-103 leído de su datasheet. El pico calculado es de unos \SI{0.28}{\ampere}; queda pendiente comparar ese pico con el dato de Coilcraft.
\end{ojo}

\subsection{Los condensadores}\label{sec:c02-cond}

\paragraph{CIN y CHF (entrada).} Cuando el interruptor de arriba se cierra, la corriente de entrada salta de cero a la corriente de la bobina en nanosegundos. Esa corriente no puede venir de la fuente de banco (el cable tiene inductancia): sale de los condensadores de entrada. TI pide \doc{SNVSB49D p.23}:
\begin{itemize}
  \item un mínimo de \SI{4.7}{\micro\farad} cerámicos, de tensión nominal al menos igual a la máxima de entrada y preferiblemente el doble. Aquí: 2$\times$\SI{4.7}{\micro\farad} de \SI{100}{\volt} por buck, más de cuatro veces los \SI{24}{\volt};
  \item además, un \SI{220}{\nano\farad} pequeño \textbf{lo más cerca posible de cada par VIN/PGND}, también de \SI{100}{\volt} y X7R. El integrado tiene VIN y PGND en los dos lados (pines 2/1 y 10/11), así que lleva dos: \net{C29} y \net{C30}.
\end{itemize}

\begin{decision}[los \SI{220}{\nano\farad} se quedan en \SI{100}{\volt}]
\textbf{Qué se decidió:} conservar el condensador original de \SI{100}{\volt} (C0805C224K1RACTU, LCSC C2167405) en \net{C29}, \net{C30}, \net{C39} y \net{C40}.\\
\textbf{Alternativas:} la auditoría de piezas propuso uno \emph{basic} de \SI{50}{\volt} (C5378), más barato en JLCPCB.\\
\textbf{Por qué:} el datasheet dice que ``the \SI{220}{\nano\farad} must also be rated at \SI{100}{\volt} with an X7R dielectric''. Con \SI{50}{\volt} se cumpliría el ``doble de la entrada'', pero no lo que pide TI para esta pieza.\\
\textbf{Fuente:} \doc{SNVSB49D p.23}; \archivo{PLAN_RUTEO.md} §6; \archivo{tools/design.py} (cambio C17).
\end{decision}

\paragraph{COUT (salida).} Los condensadores de salida filtran el rizado de la bobina y sostienen la tensión cuando la carga cambia de golpe. TI lista 3$\times$\SI{15}{\micro\farad} como opción optimizada para transitorios y 2$\times$\SI{15}{\micro\farad} como mínimo a \SI{1}{\mega\hertz} \doc{SNVSB49D p.19}. Esta placa usa 4$\times$\SI{15}{\micro\farad} de \SI{16}{\volt}, X7R, 1210 (TDK C3225X7R1C156MT000N) por buck: uno más que la opción recomendada.

\begin{idea}[la capacidad ``efectiva'' de un cerámico]
Un condensador cerámico X7R o X5R \textbf{pierde capacidad cuando tiene tensión DC encima}. Un \SI{10}{\micro\farad} de \SI{10}{\volt} en 0603 con \SI{5}{\volt} aplicados puede quedar en una fracción de su valor nominal. TI lo advierte: todos los valores de su guía son \emph{efectivos}, es decir, ya con la tensión aplicada; los encapsulados grandes y las tensiones nominales altas pierden menos \doc{SNVSB49D p.18}. Por eso en esta placa los condensadores de potencia son grandes (2220, 1210) y de tensión nominal holgada.
\end{idea}

\begin{ojo}[capacidad efectiva y CFF no cuantificados]
\begin{itemize}
  \item El cambio C17 de \archivo{tools/design.py} dice que la capacitancia efectiva de los nuevos CIN y COUT se verificó, pero el repositorio no guarda los números de las curvas de DC bias. No hay un valor efectivo citado aquí.
  \item La tabla de TI pone un condensador de \emph{feedforward} (CFF) de \SI{20}{\pico\farad} en paralelo con $R_{FBT}$. El diseño del equipo no lo tiene. El datasheet lo exige solo cuando $R_{FBT}=\SI{1}{\mega\ohm}$ \doc{SNVSB49D p.20--21} y dice que con otros valores puede mejorar el transitorio o el margen de fase ``en algunos casos'' \doc{SNVSB49D p.23}. Con \SI{100}{\kilo\ohm} y la carga casi constante de esta placa no es obligatorio, pero la respuesta a un escalón de carga no se ha medido \medir.
\end{itemize}
\end{ojo}

\paragraph{CBOOT y CVCC.} El MOSFET de arriba necesita una tensión de compuerta por encima de VIN. La saca de \textbf{CBOOT}, un condensador entre BOOT y SW que se recarga en cada ciclo: TI pide \SI{100}{\nano\farad} de al menos \SI{16}{\volt} \doc{SNVSB49D p.23}, y aquí se usa \SI{100}{\nano\farad} X7R de \SI{50}{\volt} (\net{C35}, \net{C45}). \textbf{VCC} es la salida del regulador interno de \SI{5}{\volt} del integrado, que alimenta su lógica: pide \SI{1}{\micro\farad} de \SI{16}{\volt} a tierra y ninguna carga externa \doc{SNVSB49D p.4 y p.23}. Aquí lleva \SI{2.2}{\micro\farad} X7R 0805 (\net{C36}, \net{C46}): el valor nominal mayor deja margen para la pérdida por DC bias con \SI{5}{\volt} encima.

\section{Tiempos mínimos de encendido y apagado}\label{sec:c02-tiempos}

\begin{idea}[por qué hay un pulso mínimo]
Un MOSFET integrado no puede encenderse y apagarse en un tiempo arbitrariamente corto: el driver tarda, y el integrado necesita un tiempo mínimo para medir la corriente. El datasheet da un \textbf{tiempo mínimo de encendido} $t_{ON,min}$ y un \textbf{tiempo mínimo de apagado} $t_{OFF,min}$. Si el ciclo de trabajo pide un pulso más corto que $t_{ON,min}$, el integrado no lo puede generar tal cual. Por eso hay que comprobarlo con la entrada más alta (que da el $D$ más chico) y la frecuencia más alta (que da el período más corto).
\end{idea}

Para el LMR36015: $t_{ON,min}$ de \SI{55}{\nano\second} típico y \SI{83}{\nano\second} máximo; $t_{OFF,min}$ de \SI{53}{\nano\second} típico y \SI{73}{\nano\second} máximo. La frecuencia de la variante de \SI{1}{\mega\hertz} puede estar entre \SIrange{0.85}{1.15}{\mega\hertz} \doc{SNVSB49D p.7}. Se compara contra los máximos.

\begin{calculo}[el peor caso]
$t_{ON} = \dfrac{D}{f_{SW}} = \dfrac{V_{OUT}}{V_{IN}\cdot f_{SW}}$. El peor caso es la salida de \SI{3.3}{\volt}, con \SI{28.9}{\volt} de entrada (la ruptura mínima del TVS) y el oscilador en \SI{1.15}{\mega\hertz}:
\[
t_{ON} = \frac{3{,}3}{28{,}9\cdot\SI{1.15}{\mega\hertz}} = \SI{99}{\nano\second} > \SI{83}{\nano\second}
\]
Margen: \SI{16}{\nano\second}. El apagado nunca es problema, porque con $D$ pequeño $t_{OFF}$ es largo (más de \SI{600}{\nano\second} en todos los casos).
\end{calculo}

\begin{table}[htbp]
\centering
\caption{Tiempos de encendido en los casos extremos (extracto de \archivo{CALCULOS.md} §2). Todos superan los \SI{83}{\nano\second} máximos.}
\label{tab:c02-ton}
\begin{tabular}{@{}lccrr@{}}
\toprule
\textbf{Salida} & $V_{IN}$ & $f_{SW}$ & $t_{ON}$ & \textbf{Margen}\\
\midrule
\SI{3.3}{\volt} & \SI{23.5}{\volt} (normal) & \SI{1.15}{\mega\hertz} & \SI{122}{\nano\second} & \SI{39}{\nano\second}\\
\SI{3.3}{\volt} & \SI{26.0}{\volt} ($V_{RWM}$ del TVS) & \SI{1.15}{\mega\hertz} & \SI{110}{\nano\second} & \SI{27}{\nano\second}\\
\SI{3.3}{\volt} & \SI{28.9}{\volt} ($V_{BR}$ mín.\ del TVS) & \SI{1.15}{\mega\hertz} & \SI{99}{\nano\second} & \SI{16}{\nano\second}\\
\SI{5.0}{\volt} & \SI{28.9}{\volt} & \SI{1.15}{\mega\hertz} & \SI{150}{\nano\second} & \SI{67}{\nano\second}\\
\bottomrule
\end{tabular}
\end{table}

Conclusión: en ningún caso de entrada (de \SI{18}{\volt} a \SI{28.9}{\volt}) y de frecuencia el integrado tiene que hacer un pulso más corto que su mínimo, así que \textbf{no hay repliegue de frecuencia} ni saltos de pulsos. Eso mantiene la conmutación en \SI{1}{\mega\hertz}, lejos del audio.

\section{El riel analógico 5V\_A}\label{sec:c02-5va}

El buck de \SI{5}{\volt} tiene un rizado a \SI{1}{\mega\hertz} y armónicos. Los op-amps y la parte analógica del ADC lo ``escuchan'' a través de su alimentación. Para limpiarlo, entre \net{5V} y \net{5V_A} hay un filtro pasa-bajos RC: \net{R33} de \SI{10}{\ohm} en serie y \net{C47} de \SI{10}{\micro\farad} a tierra (figura~\ref{fig:c02-rc}).

\begin{figure}[htbp]
\centering
\begin{circuitikz}
  \draw (0,2) node[left]{\net{5V} (de \net{U7})} to[short,o-] (0.4,2) to[R, l=\net{R33} \SI{10}{\ohm}] (2.8,2) to[short,-*] (3.6,2)
        to[short,-o] (5.4,2) node[right]{\net{5V_A} $\to$ \net{U3.8}, \net{U9.8}, \net{U4.3}};
  \draw (3.6,2) to[C, l=\net{C47} \SI{10}{\micro\farad}] (3.6,0) node[ground]{};
\end{circuitikz}
\caption{Filtro RC del riel analógico. Además de \net{C47}, cada carga tiene su propio desacoplo pegado al pin (\net{C7}, \net{C54}, \net{C17}, \net{C18}).}
\label{fig:c02-rc}
\end{figure}

\begin{calculo}[frecuencia de corte y atenuación]
\[
f_c = \frac{1}{2\pi R C} = \frac{1}{2\pi\cdot\SI{10}{\ohm}\cdot\SI{10}{\micro\farad}} \approx \SI{1.59}{\kilo\hertz}
\]
Un filtro RC de primer orden cae \SI{20}{\decibel} por década por encima de $f_c$. A \SI{1}{\mega\hertz}:
\[
A \approx 20\log_{10}\frac{\SI{1}{\mega\hertz}}{\SI{1.59}{\kilo\hertz}} = 20\log_{10}(629) \approx \SI{56}{\decibel}
\]
Es decir, el rizado de conmutación llega a \net{5V_A} unas 630 veces más pequeño.
\end{calculo}

\begin{figure}[htbp]
\centering
\begin{tikzpicture}
\begin{axis}[width=0.85\linewidth, height=5cm, xmode=log, xmin=10, xmax=1e7, ymin=-75, ymax=5,
  xlabel={frecuencia (Hz)}, ylabel={$|H|$ (dB)}, grid=both, grid style={gray!20},
  ytick={-70,-60,-50,-40,-30,-20,-10,0}]
  \addplot[thick, verdeplaca, domain=1:7, samples=200] ({10^x}, {-10*log10(1+(10^x/1591.5)^2)});
  \draw[dashed, cobre] (axis cs:1591.5,-75) -- (axis cs:1591.5,5);
  \node[cobre, font=\footnotesize, anchor=west, fill=white, inner sep=1pt] at (axis cs:1750,-45) {$f_c\approx\SI{1.6}{\kilo\hertz}$};
  \draw[dashed, azulidea] (axis cs:1e6,-75) -- (axis cs:1e6,5);
  \node[azulidea, font=\footnotesize, anchor=north east, fill=white, inner sep=1pt] at (axis cs:9e5,-2) {\SI{1}{\mega\hertz}: $-56$\,dB};
  \fill[azulidea!12] (axis cs:20,-75) rectangle (axis cs:20000,-68);
  \node[font=\scriptsize, azulidea] at (axis cs:650,-71.5) {banda de audio};
\end{axis}
\end{tikzpicture}
\caption{Respuesta ideal del filtro \net{R33}/\net{C47}. Atenúa fuerte la conmutación de \SI{1}{\mega\hertz}. En la banda de audio filtra poco: ahí el rechazo a la alimentación lo dan los propios op-amps y el ADC.}
\label{fig:c02-bode}
\end{figure}

\begin{ojo}[el filtro real es algo peor]
\net{C47} es un \SI{10}{\micro\farad} X5R de \SI{10}{\volt} en 0603. Con \SI{5}{\volt} encima pierde capacidad (DC bias), así que la frecuencia de corte real queda algo más arriba que \SI{1.6}{\kilo\hertz}. A \SI{1}{\mega\hertz} también pesa la inductancia parásita del condensador. La atenuación real se confirma con el osciloscopio midiendo el rizado en \net{5V_A} \medir.
\end{ojo}

\subsection*{El precio del filtro: la caída en R33}

Toda la corriente de \net{5V_A} pasa por \net{R33}:

\begin{calculo}[caída y potencia en \net{R33}]
Con la cota de \SI{14.0}{\milli\ampere}: $\Delta V = \SI{10}{\ohm}\cdot\SI{14}{\milli\ampere} = \SI{140}{\milli\volt}$ y $P = I^2R = \SI{2.0}{\milli\watt}$, muy por debajo de los \SI{100}{\milli\watt} de una resistencia 0603 (\archivo{CALCULOS.md} §4 y §6).

Si a eso le sumas la tolerancia de $\pm\SI{1.5}{\percent}$ del buck, \net{5V_A} queda entre unos \SI{4.79}{\volt} (\SI{5}{\volt} menos el \SI{1.5}{\percent}, con \SI{14}{\milli\ampere}) y \SI{4.97}{\volt} (\SI{5}{\volt} más el \SI{1.5}{\percent}, con la corriente típica de unos \SI{10.8}{\milli\ampere}). Ese rango se usa en el capítulo~\ref{cap:analog} para la región de transición del TLV9152.
\end{calculo}

\begin{decision}[un riel analógico separado]
\textbf{Qué se decidió:} alimentar los op-amps y el VCC del ADC desde \net{5V_A}, filtrado con RC, y lo digital desde \net{3V3}, con otro buck.\\
\textbf{Alternativas:} un solo riel para todo, o un regulador lineal (LDO) después del buck en lugar del RC.\\
\textbf{Por qué:} separar lo analógico de lo digital evita que las corrientes de conmutación de la lógica (relojes, I$^2$S) pasen por la alimentación de los nodos sensibles. El RC es la forma más simple de quitar el rizado de \SI{1}{\mega\hertz}: dos piezas \emph{basic} y \SI{2}{\milli\watt} de pérdida. Un LDO filtraría también parte de la banda de audio, pero agrega un integrado y su propia caída.\\
\textbf{Fuente:} \archivo{hardware/docs/connection_matrix.md} (el filtro R501/C501 del diseño original); \archivo{tools/design.py}.
\end{decision}

\section{La NUCLEO y el orden de encendido}\label{sec:c02-nucleo}

La NUCLEO-L432KC trae su propio regulador y \textbf{se alimenta por su cable USB}, el mismo que la conecta al PC. Con la placa solo comparte tierra y señales: sus pines \net{+3V3}, \net{+5V} y \net{VIN} \textbf{no se conectan} al zócalo (cambio C12). Así no hay dos fuentes de \SI{3.3}{\volt} en paralelo peleándose (la de la NUCLEO y el buck \net{U6}).

El orden recomendado es: \textbf{primero el USB, después los \SI{24}{\volt}}.

\begin{idea}[por qué importa el orden]
La razón probable es la siguiente (es un razonamiento, no está escrito en el repositorio). Las resistencias de pull-up del I$^2$C (\net{R10}, \net{R11}) van al \net{3V3} de la placa. Si la placa se enciende con la NUCLEO apagada, esas resistencias ponen \SI{3.3}{\volt} en pines de un microcontrolador sin alimentación. La corriente entra por los diodos de protección internos del pin y puede alimentar parcialmente el micro ``por la puerta de atrás''. Encender primero la NUCLEO lo evita.
\end{idea}

\section{Lo que importa para el layout de esta etapa}

Un buck tiene un \textbf{lazo caliente}: el camino CIN/CHF $\to$ VIN $\to$ MOSFET $\to$ PGND $\to$ CIN, por donde la corriente cambia en nanosegundos. Cuanto más área encierra ese lazo, más inductancia tiene y más ruido irradia. Por eso los \SI{220}{\nano\farad} van pegados a cada par VIN/PGND y todo el lazo de potencia queda en la capa superior. El capítulo~\ref{cap:ruteo} trata este layout a fondo, con las mediciones de los lazos.

\section*{Resumen}
\begin{itemize}
  \item La placa consume poco: \SI{22.4}{\milli\ampere} en \net{3V3}, \SI{14.0}{\milli\ampere} en \net{5V_A} y unos \SI{25}{\milli\ampere} de la fuente de \SI{24}{\volt} \estimado.
  \item \net{D1} (Schottky en serie) protege contra la polaridad invertida y cuesta \SI{0.5}{\volt} y \SI{20}{\milli\watt}. \net{D2} (TVS de \SI{26}{\volt}) recorta los picos.
  \item Un buck trocea la tensión con un ciclo de trabajo $D=V_{OUT}/V_{IN}$. Es mucho más eficiente que un lineal, que aquí perdería unos \SI{0.45}{\watt} solo en el riel de \SI{3.3}{\volt}.
  \item El LMR36015FB conmuta siempre a \SI{1}{\mega\hertz} (FPWM): el ruido nunca cae en la banda de audio, a cambio de hasta unos \SI{0.15}{\watt} de pérdida en vacío por integrado \estimado.
  \item Los divisores \SI{100}{\kilo\ohm}/\SI{43.2}{\kilo\ohm} y \SI{100}{\kilo\ohm}/\SI{24.9}{\kilo\ohm} dan \SI{3.315}{\volt} y \SI{5.016}{\volt}. Las bobinas dan un rizado del 26--28\,\% de \SI{1.5}{\ampere}, como recomienda TI.
  \item El pulso más corto (\SI{99}{\nano\second}) supera el mínimo de \SI{83}{\nano\second}: no hay repliegue de frecuencia.
  \item \net{R33}/\net{C47} filtran \net{5V} con $f_c\approx\SI{1.6}{\kilo\hertz}$: unos \SI{56}{\decibel} a \SI{1}{\mega\hertz}, a cambio de \SI{140}{\milli\volt} de caída.
  \item La NUCLEO se alimenta por su USB; se enciende antes que los \SI{24}{\volt}.
\end{itemize}
